\section{Related Work}
\label{sec:related}

\paragraph{Reference counting and cycles.}
Counting~\cite{collins1960} is the oldest automatic reclamation and the one closest to this
design's barrier, but its classical costs are a write barrier on every store and cycles that no
count can break \cite{jones2011handbook}. Cycle collection has been attacked with trial deletion
and concurrent variants \cite{bacon2001concurrent}, with deferred and hybrid counting
\cite{blackburn2003ulterior}, and with reuse-based optimizations in a functional setting
\cite{reinking2021perceus}.  extC keeps a counting mode but makes it a tier that the compiler
chooses, not the default: the default emits no counter at all, and the shapes where counting would
pay — a reference stored into a long-lived slot and dropped later — are covered by a budgeted sweep
whose trigger is memory pressure rather than a decrement to zero.

\paragraph{Tracing collectors and bounded pause.}
Tracing through roots \cite{jones2011handbook} can move objects and can schedule work on other
threads, and real-time variants bound pauses by work-based scheduling \cite{bacon2003realtime}.
Uncooperative environments require conservative root discovery \cite{boehm1988uncooperative}, and
mark-region collectors combine tracing with size-class allocation \cite{blackburn2008immix}, which
is the closest published relative of our sweep tier.  The difference we exploit is that extC emits
the root set statically: there is no stack scan, so the collector never needs to interpret a frame,
and a function that provably stores no escaping reference gets no code from the runtime at all.

\paragraph{Regions, ownership, and lifetimes.}
Region inference \cite{tofte1997region} and region-based dialects of C
\cite{grossman2002cyclone,gay2001regions} place objects in lexically scoped regions and free a
region as a unit, which is what our frame arena and the wholesale release of \S\ref{sec:design} do.
Ownership and linear type systems make deallocation explicit and checkable
\cite{clarke1998ownership,bernardy2018linearhaskell}, Rust's borrow checker is the industrial
form of the same idea \cite{jung2018rustbelt}, and reference capabilities with actor-local heaps
show how far the discipline can be pushed concurrently \cite{clebsch2013pony}.
extC keeps the region half of this lineage and not the ownership half: values are copied, mutable
references may alias, and the checker proves only that a view does not outlive its storage.  That
weaker discipline is what allows the fallback ladder: when the static analysis cannot prove a free,
the compiler emits the reason and the runtime finishes the job under a budget.

\paragraph{Escape analysis and placement.}
Escape analysis decides where an allocation may live \cite{choi1999escape,whaley1999rinard}, and
partial escape analysis with scalar replacement can remove allocations entirely
\cite{stadler2014partial}.  In extC escape depth is not only a placement decision: it is the number
that the root table, the cross-region reference bit and the acyclicity judgement are all stated
in.  Promotion is therefore a static decision — the allocation is born in the shallow region —
rather than a runtime copy, which is why the safety proof of \S\ref{sec:formal} needs no move rule.

\paragraph{Analytic performance models.}
The roofline model \cite{williams2009roofline} puts a workload on a plane of achieved rate against
operational intensity and says which of two hardware ceilings binds.  \S\ref{sec:roofline} makes
the same move for reclamation policy, with allocation bytes as the unit of work, collector time per
allocated byte as the intensity, and the memory and pause contracts as the two ceilings.  The
analogue is not a metaphor: the cost law, the memory law, the pause law and the feasibility
condition are theorems about the same measured constants, and the policy question ("which batch and
which trigger ratio") is a corner of a two-dimensional region, as it is on a roofline.

\paragraph{The cost of automatic memory management.}
Hertz and Berger \cite{hertz2005quantifying} measured collectors against explicit management and
found the trade to be roughly a factor of two in time for a factor of three to four in memory; our
numbers have the same shape but invert the sign on allocation-bound programs, because the
alternative we compare against is not a tuned allocator but a bump arena that never reuses memory
(\tab\ref{tab:ablate}: 2.1 to 3.7 times the time and 19 to 59 times the memory with reclamation
off).  The analogy has one boundary worth naming: the store barrier is the cost that cannot be
removed when a reference is stored, and it is where the pointer-chasing shape trades a few percent
of time for 13\% less peak memory than the best C++ arm.

\paragraph{What is new here, stated narrowly.}
Each ingredient has a precedent: block-exit region release \cite{tofte1997region}, static
placement by escape depth \cite{choi1999escape}, incremental mark--sweep with a work budget
\cite{blackburn2008immix,bacon2003realtime}, and reclamation at the store that orphans a structure
\cite{collins1960}.  The contribution is three consequences of putting them in one compiler and
asking, per unit of compilation, whether each is needed at all:
\begin{enumerate}
\item the root set is a static data structure emitted from checking, so there is no stack scan and a
unit that needs nothing gets no runtime code — 68.5\% of the corpus programs that compile emit
none (\S\ref{sec:eval-census});
\item the tiers can be turned off independently, and the no-reclamation build is a working compiler
mode rather than a strawman, which is what makes the mechanism ablation of
\S\ref{sec:eval-ablation} comparable arm by arm;
\item the resulting machine is small enough to state contracts about: seven theorems, four
calibrated constants, a Lean development of the region and idle invariants, and a gate suite that
fails when the emitted text stops matching the theorems.
\end{enumerate}

\paragraph{Alternatives we did not measure.}
Region-based systems with reference counting for escaping values \cite{grossman2002cyclone},
Pony-style reference capabilities \cite{clebsch2013pony}, and ownership-passing variants of
Perceus occupy adjacent points; we compare against hand-written C++, Rust and Go because those are
what a practitioner would reach for, not because they are the closest research systems.  A
Perceus-style reuse analysis on these workloads is the obvious missing experiment, and a
concurrent version of our contract is the obvious missing theory.
